% Part: sets-functions-relations
% Chapter: sets
% Section: important-sets
% Hindi translation (hi-Deva-IN) of Open Logic commit 9620cc73f9c8e0ad003c514a5d3748f29611c4c0.

\documentclass[../../../include/open-logic-section]{subfiles}

\begin{document}

\olfileid[hi]{sfr}{set}{imp}
\olsection{कुछ महत्त्वपूर्ण समुच्चय}

\begin{ex}
हम मुख्यतः ऐसे समुच्चयों पर विचार करेंगे जिनके !!{element} गणितीय वस्तुएँ
हैं। ऐसे चार समुच्चय इतने महत्त्वपूर्ण हैं कि उनके विशिष्ट नाम हैं:
\begin{multline*}
    \Nat = \{0, 1, 2, 3, \ldots\} \\
    \shoveright{\text{प्राकृत संख्याओं का समुच्चय}}\\
    \shoveleft{\Int = \{\ldots, -2, -1, 0, 1, 2, \ldots\}} \\
    \shoveright{\text{पूर्णांकों का समुच्चय}}\\
    \shoveleft{\Rat = \Setabs{\nicefrac{m}{n}}{m, n \in \Int\text{ और }n \neq 0}}\\
    \shoveright{\text{परिमेय संख्याओं का समुच्चय}}\\
    \shoveleft{\Real = (-\infty, \infty)}\\
    \text{वास्तविक संख्याओं का समुच्चय (सातत्य)}
\end{multline*}
ये सभी \emph{अनंत} समुच्चय हैं, अर्थात् इनमें से प्रत्येक में अनंत संख्या में
!!{element} हैं।

इन समुच्चयों में क्रमशः आगे बढ़ते हुए हम अपने संख्या-संग्रह में \emph{और}
संख्याएँ जोड़ते जाते हैं। वास्तव में यह स्पष्ट होना चाहिए कि
$\Nat \subseteq \Int
\subseteq \Rat \subseteq \Real$: आखिरकार, प्रत्येक
प्राकृत संख्या पूर्णांक है; प्रत्येक पूर्णांक परिमेय है; और प्रत्येक परिमेय
संख्या वास्तविक है। इसी प्रकार यह स्पष्ट होना चाहिए कि
$\Nat \subsetneq \Int \subsetneq
\Rat$, क्योंकि $-1$ पूर्णांक तो है किंतु
प्राकृत संख्या नहीं, और $\nicefrac{1}{2}$ परिमेय तो है किंतु पूर्णांक नहीं।
यह कम स्पष्ट है कि $\Rat \subsetneq \Real$, अर्थात् कुछ वास्तविक संख्याएँ
परिमेय नहीं हैं\oliflabeldef{sfr:arith:real:realline}{, किंतु हम
\olref[arith][real]{realline} में इस विषय पर लौटेंगे}{}।

हम कभी-कभी धन पूर्णांकों के समुच्चय $\PosInt = \{1,
2, 3, \dots\}$ और केवल
पहली दो प्राकृत संख्याओं वाले समुच्चय $\Bin = \{0, 1\}$ का भी उपयोग करेंगे।
\end{ex}

\begin{tagblock}{compsci}
\begin{ex}[वर्ण-शृंखलाएँ]
एक अन्य रोचक उदाहरण समुच्चय~$A^{*}$ है। यह वर्णमाला~$A$ पर \emph{परिमित
वर्ण-शृंखलाओं} का समुच्चय है: $A$ के अवयवों का कोई भी परिमित अनुक्रम, $A$
पर एक वर्ण-शृंखला है। हम \emph{रिक्त वर्ण-शृंखला $\Lambda$} को $A$ पर बनी
वर्ण-शृंखलाओं में भी सम्मिलित करते हैं, और ऐसा प्रत्येक वर्णमाला~$A$ के लिए
करते हैं। उदाहरणतः,
\begin{multline*}
\Bin^*
=\{\Lambda,0,1,00,01,10,11,\\
000,001,010,011,100,101,110,111,0000,\ldots\}.
\end{multline*}
यदि $x=x_{1}\ldots x_{n}\in A^{*}$, $n$ ``प्रतीकों'' से बनी $A$ पर एक
वर्ण-शृंखला है, तो हम कहते हैं कि उस वर्ण-शृंखला की \emph{लंबाई} $n$ है और
$\len{x}=n$ लिखते हैं।
\end{ex}
\end{tagblock}

\begin{ex}[अनंत अनुक्रम]
किसी भी समुच्चय~$A$ के लिए हम समुच्चय~$A^\omega$ पर भी विचार कर सकते हैं,
जो $A$ के !!{element}ों के अनंत अनुक्रमों का समुच्चय है। अनंत अनुक्रम
$a_1a_2a_3a_4\dots$ वस्तुओं की एकदिश अनंत सूची से बना होता है, जिसकी प्रत्येक
वस्तु $A$ का !!a{element} होती है।
\end{ex}

\end{document}
